\documentclass[12pt]{article}
\usepackage[utf8]{inputenc}
\usepackage[T1]{fontenc}
\usepackage{hyperref}
\title{Highland Totonac Language Resources for K-12 Curriculum Development}
\author{Compiled by Claude Sonnet 5 (claude-sonnet-5) for Andr\'as Kornai}
\date{}
\begin{document}
\maketitle

\section{Introduction}
Highland (Sierra) Totonac (ISO 639-3: tos; Glottolog: high1243) is a Totonacan language of Mexico. Searches found few dedicated teaching materials. What turned up were catalogue entries and aggregator pages, so this report is short.

\section{Primary Educational Resources}
The Endangered Languages Project page for Sierra Totonac lists a grammar-category resource, \emph{Aprendamos totonaco} (2010). I could not confirm its author, publisher or whether it is a classroom text.
\begin{itemize}
\item Endangered Languages Project, Sierra Totonac resources: \url{https://endangeredlanguages.com/fil/node/38946/resources}
\end{itemize}

\section{Linguistic and Reference Materials}
\begin{itemize}
\item Glottolog entry for Highland Totonac. It cross-references descriptive literature, but the search results did not show specific titles: \url{https://glottolog.org/resource/languoid/id/high1243}
\item Glosbe dictionary pages for Highland Totonac. These are aggregated and not a vetted scholarly dictionary: \url{https://glosbe.com/tos/de}
\end{itemize}

\section{Religious and Community Resources}
No religious or community-organisation resources for tos appeared in the search results.

\section{Digital and Audio Resources}
No audio or digital teaching resources were found beyond the aggregator pages above.

\section{Curriculum Development Context}
I found no evidence of school programmes in Highland Totonac. A curriculum would likely need new materials. The first step would be locating the \emph{Aprendamos totonaco} text and scholarly grammars through library catalogues, which these searches did not reach.
\end{document}